\documentclass[12pt,a4paper]{article}
\usepackage[utf8]{inputenc}
\usepackage[T1]{fontenc}
\usepackage[spanish,es-noshorthands]{babel}
\usepackage{geometry}
\geometry{a4paper, margin=2.5cm}
\usepackage{booktabs}
\usepackage{graphicx}
\usepackage{hyperref}
\usepackage{xcolor}
\usepackage{tabularx}
\usepackage{titling}

\hypersetup{
    colorlinks=true,
    linkcolor=blue,
    filecolor=magenta,      
    urlcolor=blue,
    pdftitle={Documentación ESP32 Kids Lab},
}

\title{\textbf{Documentación Técnica: ESP32 Kids Lab}}
\author{Proyecto: Laboratorio ESP32}
\date{\today}

\begin{document}

\maketitle

\section{Descripción General}
\textbf{ESP32 Kids Lab} es una plataforma educativa interactiva diseñada para que los niños exploren conceptos básicos de electrónica, sensores y actuadores. El proyecto convierte un ESP32 en un laboratorio de bolsillo autónomo, exponiendo una red Wi-Fi propia (SoftAP) y un servidor web interactivo con un diseño moderno y amigable.

\section{Características Principales}
\begin{itemize}
    \item \textbf{Conexión Directa:} El ESP32 crea su propia red Wi-Fi (\texttt{Laboratorio ESP32}) por lo que no requiere internet ni un router externo.
    \item \textbf{Interfaz Web Moderna:} Creada con HTML5, CSS3 y JavaScript puro. Se sirve directamente desde la memoria flash del ESP32 usando \textbf{LittleFS}.
    \item \textbf{Sensores en Vivo:} Lectura en tiempo real de Distancia, Luz y Temperatura.
    \item \textbf{Control de Actuadores:} Interruptores en la web para encender y apagar LEDs de colores, un Relé y un Buzzer.
    \item \textbf{Taller Mágico (Display 7 Segmentos):} Resuelve sumas matemáticas desde la web y el resultado aparece mágicamente en el display físico TM1637.
    \item \textbf{Estudio Pixel Art (Matriz NeoPixel):} Una cuadrícula interactiva de 4x4 en la web que permite pintar y dibujar en tiempo real sobre una matriz física de LEDs WS2812.
\end{itemize}

\section{Mapa de Pines (Hardware Pinout)}

A continuación se detalla la asignación de pines (GPIO) del ESP32 utilizados en el proyecto:

\vspace{0.5cm}
\noindent
\begin{table}[h]
\centering
\begin{tabularx}{\textwidth}{llcX}
\toprule
\textbf{Componente} & \textbf{Tipo} & \textbf{Pin (GPIO)} & \textbf{Notas} \\
\midrule
\textbf{LED Rojo} & Actuador / Salida & \texttt{1} (TX) & Su uso desactiva la salida del Monitor Serial. \\
\textbf{LED Verde} & Actuador / Salida & \texttt{23} & \\
\textbf{LED Azul} & Actuador / Salida & \texttt{19} & \\
\textbf{LED Amarillo} & Actuador / Salida & \texttt{3} (RX) & \\
\textbf{Relé} & Actuador / Salida & \texttt{0} & Control de cargas (escuchar el ``clic''). \\
\textbf{Buzzer} & Actuador / Salida & \texttt{12} & Requiere activar el Switch 1 en la placa. \\
\textbf{Trigger} & Sensor / Salida & \texttt{5} & Pin de disparo ultrasónico (HC-SR04). \\
\textbf{Echo} & Sensor / Entrada & \texttt{18} & Pin de eco ultrasónico (HC-SR04). \\
\textbf{LDR} & Sensor / Entrada Analógica & \texttt{34} & Fotorresistencia (medición de Luz). \\
\textbf{LM35} & Sensor / Entrada Analógica & \texttt{35} & Sensor de temperatura. \\
\textbf{CLK} & Actuador / Reloj & \texttt{27} & Reloj para el display TM1637. \\
\textbf{DIO} & Actuador / Datos & \texttt{14} & Datos para el display TM1637. \\
\textbf{Matriz NeoPixel} & Actuador / Datos & \texttt{13} & Control para la Matriz WS2812 de 16 LEDs. \\
\bottomrule
\end{tabularx}
\end{table}

\section{Arquitectura de Software}

El proyecto se divide claramente en Backend (C++ en ESP32) y Frontend (Web en LittleFS):

\subsection{Backend (ESP32 - PlatformIO)}
\begin{itemize}
    \item \textbf{Framework:} Arduino Core para ESP32.
    \item \textbf{ESPAsyncWebServer:} Maneja las peticiones HTTP de forma asíncrona, permitiendo servir la interfaz web y responder a las APIs REST sin bloquear el ciclo principal del procesador.
    \item \textbf{API REST (\texttt{main.cpp}):}
    \begin{itemize}
        \item \texttt{GET /api/sensors}: Devuelve un objeto JSON con los valores actuales de distancia, luz y temperatura.
        \item \texttt{POST /api/actuators}: Recibe un JSON booleano para encender/apagar de forma independiente los LEDs, el Relé y el Buzzer.
        \item \texttt{POST /api/display}: Recibe un número para mostrarlo directamente en el módulo TM1637.
        \item \texttt{POST /api/matrix}: Interfaz avanzada para la matriz NeoPixel. Recibe comandos para limpiar, pintar todo de un color, o un array complejo. Usa un \texttt{DynamicJsonDocument} de 2048 bytes para soportar payloads grandes.
    \end{itemize}
\end{itemize}

\subsection{Frontend (LittleFS)}
\begin{itemize}
    \item \textbf{\texttt{index.html}:} Estructura semántica dividida en tarjetas: Paneles de Sensores, Actuadores, Taller Mágico, y Pixel Art.
    \item \textbf{\texttt{style.css}:} Sistema de diseño basado en variables CSS. Utiliza flexbox, grillas y animaciones suaves para generar una experiencia moderna inspirada en el \textit{Glassmorphism}.
    \item \textbf{\texttt{app.js}:} 
    \begin{itemize}
        \item Ciclo de \textit{polling} automático (cada 2 segundos) consultando a \texttt{GET /api/sensors}.
        \item Funciones de envío asíncrono (\textit{Fetch API}) para controlar cada dispositivo interactivo.
        \item Lógica de dibujo Pixel Art, manejo de array de colores en tiempo real y técnica de \textit{debounce}.
    \end{itemize}
\end{itemize}

\section{Dependencias (Libraries)}
Las bibliotecas instaladas en el entorno (\texttt{platformio.ini}) son:
\begin{itemize}
    \item \texttt{me-no-dev/ESP Async WebServer}: Para el servidor web asíncrono y los Endpoints REST.
    \item \texttt{bblanchon/ArduinoJson} (v6): Para serialización y deserialización rápida de datos JSON.
    \item \texttt{smougenot/TM1637}: Para el manejo del display de 4 dígitos.
    \item \texttt{adafruit/Adafruit NeoPixel}: Para el protocolo de datos de la matriz WS2812.
\end{itemize}

\end{document}
